% ======================================================================
% Chapter: Symmetric Orbifolds
% ======================================================================

\chapter{\texorpdfstring{The symmetric orbifold $\operatorname{Sym}^N(X)$}{The symmetric orbifold Sym\^{}N(X)}}\label{chap:symmetric-orbifold}
% Regime III: non-perturbative (Convention~\ref{conv:regime-tags}).

The notation $\operatorname{Sym}^N(X)$ carries three different
objects. The fixed-point chiral algebra is
\[
 X_N^{\mathrm{id}} := (X^{\otimes N})^{S_N};
\]
its ordered bar complex is built from the invariant chiral product.
The full permutation-orbifold conformal field theory has state space
\[
 \mathcal H_N^{\mathrm{orb}}
 = \bigoplus_{[\sigma]\in \operatorname{Conj}(S_N)}
 \mathcal H_\sigma^{C(\sigma)} ,
\]
with $C(\sigma)$ the centralizer and with twisted sectors indexed by
conjugacy classes. A
$G$-crossed orbifold completion, when constructed, adjoins twist
fields as genuine fields and has a different bar complex. The
genus-$1$ scalar identity
\[
 \kappa^{\mathrm{id}}_1(X_N^{\mathrm{id}})
 = N\,\kappa(X)
\]
belongs only to the first object: the identity-sector fixed-point
chiral algebra. The DMVV product \cite{DMVV97} computes the
genus-$1$ elliptic genera of the second object; it does not compute
ordered chiral bar cohomology.

The diagonal fields
$a_\Delta=\sum_{i=1}^N a^{(i)}\in (X^{\otimes N})^{S_N}$ retain the
pole orders of the seed OPE, while scalar residues add over the
$N$ independent copies. Twisted sectors have conformal weights
\[
 h_{\mathrm{twist}}(\lambda)
 = \frac{c(X)}{24}\sum_j m_j\Bigl(j-\frac1j\Bigr)
\]
for cycle type $\lambda=(1^{m_1}2^{m_2}\cdots)$, and they affect the
full orbifold partition function. They are not generators of the
ordered bar complex of $X_N^{\mathrm{id}}$.

\begin{table}[ht]
\centering
\small
\caption{Five-theorem scope for
$\operatorname{Sym}^N(X)$.}\label{tab:symn-five-theorems}
\begin{tabular}{lll}
\toprule
\textbf{Theorem} & \textbf{Statement for $\operatorname{Sym}^N(X)$}
 & \textbf{Scope} \\
\midrule
A (adjunction) &
 fixed-point dual vs.\ twisted completion
 & Conj.~\ref{conj:symn-koszul-dual} \\
B (inversion) &
 $\Omega(\barB(X_N^{\mathrm{id}}))\to X_N^{\mathrm{id}}$
 & Reynolds descent hypothesis; Prop.~\ref{conj:symn-inversion} \\
C (complementarity) &
 identity sector: $\kappa+\kappa'=N(\kappa(X)+\kappa(X'))$
 & Rem.~\ref{rem:symn-complementarity} \\
D (modular char.) &
 $\kappa^{\mathrm{id}}_1(X_N^{\mathrm{id}})=N\kappa(X)$
 & Prop.~\ref{prop:symn-kappa} \\
H (chiral Hochschild) &
 $\operatorname{ChirHoch}^n(X_N^{\mathrm{id}})=0$ for $n>2$
 & conditional on Koszul descent; Prop.~\ref{conj:symn-hochschild} \\
\bottomrule
\end{tabular}
\end{table}

\begin{table}[ht]
\centering
\small
\caption{Shadow archetype data for the fixed-point
sector.}\label{tab:symn-shadow-archetype}
\begin{tabular}{ll}
\toprule
\textbf{Invariant} & \textbf{Value} \\
\midrule
Diagonal class & Same as seed $X$ \\
Shadow depth $r_{\max}$ &
 $\geq r_{\max}(X)$ in the diagonal fixed-point channel \\
$\kappa^{\mathrm{id}}_1(X_N^{\mathrm{id}})$ &
 $N \cdot \kappa(X)$ \\
Identity-sector $F_1$ & $N \cdot \kappa(X) / 24$ \\
Koszul dual & fixed-point descent plus twisted completion data \\
Complementarity (identity sector) &
 $N \cdot (\kappa(X) + \kappa(X^!))$ \\
\bottomrule
\end{tabular}
\end{table}

\begin{proposition}[Identity-sector modular characteristic;
\ClaimStatusProvedHere]\label{prop:symn-kappa}
\index{symmetric orbifold!modular characteristic}
Let $N\geq1$, and let $X$ be a modular Koszul chiral algebra with
genus-$1$ scalar obstruction
$\mathrm{obs}_1(X)=\kappa(X)\lambda_1$. For the fixed-point chiral
algebra
$X_N^{\mathrm{id}}=(X^{\otimes N})^{S_N}$,
\[
 \kappa^{\mathrm{id}}_1(X_N^{\mathrm{id}})
 \;=\; N \cdot \kappa(X).
\]
This is the identity-sector scalar obstruction. It is not a
statement about the full twisted orbifold Hilbert space, a
higher-genus free energy, or the ordered bar complex of a
$G$-crossed orbifold completion.
\end{proposition}

\begin{proof}
The tensor product $X^{\otimes N}$ has diagonal stress tensor and
diagonal obstruction cocycle equal to the sum of the $N$ copywise
cocycles. Additivity of the scalar obstruction for independent
tensor factors gives
\[
 \mathrm{obs}_1(X^{\otimes N})
 = \sum_{i=1}^N \mathrm{obs}_1(X^{(i)})
 = N\,\kappa(X)\lambda_1 .
\]
The Reynolds idempotent
$\frac1{N!}\sum_{\sigma\in S_N}\sigma$ fixes the diagonal cocycle.
Thus the invariant chiral algebra $X_N^{\mathrm{id}}$ carries the
same diagonal scalar obstruction. No twisted-sector state is used in
this computation.
\end{proof}

\begin{remark}[K3 elliptic genus and the identity scalar]
\label{rem:symn-k3-elliptic-identity}
For the K3 sigma-model lane used here, $\kappa(K3)=2$ and
Proposition~\ref{prop:symn-kappa} gives
$\kappa^{\mathrm{id}}_1((K3^{\otimes N})^{S_N})=2N$.
The DMVV formula instead concerns the elliptic genus:
if
\[
 \operatorname{Ell}_{K3}(\tau,z)
 =\sum_{n,\ell}c_{K3}(n,\ell)q^n y^\ell,
\]
then
\[
 \sum_{N\geq0}p^N
 \operatorname{Ell}_{\operatorname{Sym}^N(K3)}(\tau,z)
 =
 \prod_{\substack{m\geq1,\; n\geq0\\ \ell}}
 (1-p^m q^n y^\ell)^{-c_{K3}(mn,\ell)} .
\]
At $z=0$ this specializes to
$\sum_N p^N\chi(\operatorname{Hilb}^N(K3))
=\prod_{m\geq1}(1-p^m)^{-24}$.
This partition-function identity includes all commuting-pair
orbifold sectors. It is compatible with the identity-sector scalar,
but it is not a proof of ordered-bar linearity.
\end{remark}

% ======================================================================
% Section: Twist sector analysis
% ======================================================================

\section{Twist sector analysis}\label{sec:symn-twist-sectors}

The full orbifold Hilbert space $\mathcal H_N^{\mathrm{orb}}$
decomposes over conjugacy classes of~$S_N$. Each conjugacy class is
parametrized by
a partition $\lambda = (1^{m_1}\, 2^{m_2} \cdots k^{m_k})$ of~$N$,
and the twist field associated to a permutation of cycle type~$\lambda$
has conformal weight
\begin{equation}\label{eq:symn-twist-weight}
 h_{\mathrm{twist}}(\lambda)
 \;=\; \frac{c(X)}{24}\, \sum_{j=1}^{k} m_j\Bigl(j - \frac{1}{j}\Bigr).
\end{equation}
The identity conjugacy class ($\lambda = (1^N)$) has
$h_{\mathrm{twist}} = 0$; after the centralizer projection it gives
the fixed-point sector $X_N^{\mathrm{id}}$. For a single
transposition ($\lambda = (2, 1^{N-2})$):
$h_{\mathrm{twist}} = c(X)/16$, the lightest twist field when
$c(X)>0$.

\begin{proposition}[Twist weights and the identity vacuum;
\ClaimStatusProvedHere]\label{prop:symn-twist-vanishing}
\index{symmetric orbifold!twist sector weights}
Assume the seed conformal field theory has $c(X)>0$. For every
nontrivial cycle type $\lambda\neq (1^N)$, the twist ground state
has $h_{\mathrm{twist}}(\lambda)>0$. Consequently no nontrivial
twist ground state has the conformal weight of the identity vacuum.
The identity-sector scalar in Proposition~\ref{prop:symn-kappa} is
therefore unaffected by adjoining the twisted Hilbert spaces.

The proposition does not assert that twisted sectors vanish from the
genus-$1$ partition function. They are precisely the additional
terms organized by the DMVV product.
\end{proposition}

\begin{proof}
For $\lambda=(1^{m_1}2^{m_2}\cdots k^{m_k})$,
\[
 h_{\mathrm{twist}}(\lambda)
 = \frac{c(X)}{24}\sum_{j=1}^k m_j\Bigl(j-\frac1j\Bigr).
\]
Each summand is nonnegative and it is positive for every cycle of
length $j>1$. A nontrivial conjugacy class contains such a cycle, so
$h_{\mathrm{twist}}(\lambda)>0$ when $c(X)>0$. The identity vacuum is
the unique state of twist weight zero. At $c(X)=0$ this positivity
argument disappears; no conclusion about the twisted contribution
follows from the weight formula alone.
\end{proof}

\begin{remark}[Twisted sectors and bar degree]
\label{rem:symn-twist-higher-degree}
\index{symmetric orbifold!twisted sector shadows}
There is no universal ``first quartic correction'' theorem. In the
fixed-point chiral algebra $X_N^{\mathrm{id}}$, twist fields are
modules rather than bar generators. In a $G$-crossed completion that
adjoins twist fields, the binary OPE
\[
 \sigma_g(z)\sigma_h(w)\sim
 (z-w)^{-\alpha(g,h)}\,\sigma_{gh}(w)+\cdots
\]
already enters the bar differential. For $N=2$ the covering-space
calculation of Lunin--Mathur \cite{LuninMathur2001} gives the leading
branch exponent $\alpha=c(X)/8$ in the $\sigma_2$--$\sigma_2$
channel. The degree at which twisted data first changes a shadow
coefficient is therefore part of the chosen orbifold completion, not
a consequence of DMVV.
\end{remark}

\begin{conjecture}[Orbifold Koszul duality;
\ClaimStatusConjectured]\label{conj:symn-koszul-dual}
\index{symmetric orbifold!Koszul dual}
The Koszul dual of the fixed-point algebra $X_N^{\mathrm{id}}$ is
obtained from the $S_N$-equivariant Koszul dual of $X^{\otimes N}$
only after the Reynolds bar comparison is a quasi-isomorphism. For a
$G$-crossed orbifold completion, additional summands are controlled
by the twist-field OPE algebra, not by $\operatorname{Sym}^N(X^!)$
alone.
\end{conjecture}

\begin{proposition}[Bar-cobar inversion under Reynolds descent;
\ClaimStatusConditional]\label{conj:symn-inversion}
The bar-cobar counit
$\Omega(\barB(X_N^{\mathrm{id}})) \to X_N^{\mathrm{id}}$
is a quasi-isomorphism if $X$ is chirally Koszul and the natural
Reynolds comparison
\[
 \barB(X_N^{\mathrm{id}})
 \longrightarrow
 \barB(X^{\otimes N})^{S_N}
\]
is a quasi-isomorphism of filtered bar complexes.
\end{proposition}

\begin{proof}
Koszulness of $X$ gives bar-cobar inversion for each tensor factor,
and Proposition~\ref{prop:independent-sum-factorization} gives the
tensor-product statement for $X^{\otimes N}$. Over characteristic
zero the finite-group invariants functor is exact on vector spaces,
so inversion descends along the displayed comparison once that
comparison is a quasi-isomorphism. The comparison is the only extra
hypothesis; it is exactly the assertion that taking the fixed-point
chiral algebra introduces no higher bar torsion.
\end{proof}

\begin{remark}[Identity-sector complementarity]
\label{rem:symn-complementarity}
\index{symmetric orbifold!complementarity}
In the identity sector, complementarity
(Theorem~\ref{thm:quantum-complementarity}) for the tensor product
$X^{\otimes N}$ gives
$\kappa(X^{\otimes N}) + \kappa((X^!)^{\otimes N}) =
N(\kappa(X) + \kappa(X^!))$. The fixed-point comparison requires the
same Reynolds descent as Proposition~\ref{conj:symn-inversion}. A
twisted orbifold completion has its own centre and its own
Hochschild complex; Theorem~C for that object is a separate theorem.
\end{remark}

\begin{proposition}[Theorem H for the fixed-point sector;
\ClaimStatusConditional]\label{conj:symn-hochschild}
If $X_N^{\mathrm{id}}$ is chirally Koszul and the compact-generation
hypotheses of Theorem~\ref{thm:hochschild-polynomial-growth} hold for
the fixed-point chiral algebra, then
\[
 \operatorname{ChirHoch}^n(X_N^{\mathrm{id}})=0
 \qquad (n>2).
\]
This is an ambient-curve statement. It does not compute the
categorical Hochschild cohomology of the full twisted orbifold
module category.
\end{proposition}

\begin{proof}
Theorem~\ref{thm:hochschild-polynomial-growth} is applied to the
single chiral algebra $X_N^{\mathrm{id}}$ on a smooth projective
curve. Its amplitude bound is the curve bound $[0,2]$. Twisted
sectors are modules for $X_N^{\mathrm{id}}$ unless a separate
$G$-crossed algebra has been constructed, so they are not inputs to
this chiral Hochschild complex.
\end{proof}


% ======================================================================
% Section: Shadow depth inheritance
% ======================================================================

\section{Shadow depth inheritance}\label{sec:symn-shadow-depth}

\begin{proposition}[Diagonal shadow depth of the fixed-point sector;
\ClaimStatusProvedHere]\label{prop:symn-shadow-depth}
\index{symmetric orbifold!shadow depth}
Let $N\geq1$, and let $X$ be a modular Koszul algebra. For the
fixed-point chiral algebra
$X_N^{\mathrm{id}}=(X^{\otimes N})^{S_N}$, the diagonal fixed-point
channels satisfy
\begin{equation}\label{eq:symn-shadow-depth-bound}
 r_{\max}(X_N^{\mathrm{id}})
 \;\geq\; r_{\max}(X).
\end{equation}
If $X$ is class~M, the diagonal Virasoro field
$T_\Delta=\sum_iT^{(i)}$ gives
$r_{\max}(X_N^{\mathrm{id}})=\infty$.
\end{proposition}

\begin{proof}
Let $u(z)v(w)$ be a seed OPE channel witnessing shadow depth~$r$.
Set $u_\Delta=\sum_i u^{(i)}$ and
$v_\Delta=\sum_i v^{(i)}$. Mixed-copy OPEs are regular, hence the
singular part of $u_\Delta(z)v_\Delta(w)$ is the sum of the $N$
copywise singular parts. The highest nonzero pole or shadow
coefficient is therefore multiplied by $N$, not cancelled. Since
the diagonal fields are $S_N$-invariant, the same channel lies in
$X_N^{\mathrm{id}}$. For class~M the seed Virasoro channel has
unbounded shadow depth, and the diagonal Virasoro channel has the
same unbounded depth.
\end{proof}

\begin{remark}[Shadow class of specific orbifolds]
\label{rem:symn-shadow-class-examples}
\begin{center}
\renewcommand{\arraystretch}{1.3}
\begin{tabular}{lccl}
\toprule
Seed $X$ & Class of $X$ & $r_{\max}(X)$ &
 Class of $\operatorname{Sym}^N(X)$ \\
\midrule
Heisenberg $\mathcal{H}_\kappa$ & G & $2$ & $\geq$ G \\
$\widehat{\mathfrak{sl}}_2$ at level $k$ & L & $3$ &
 $\geq$ L \\
$\beta\gamma$ & C & $4$ & $\geq$ C \\
$\mathrm{Vir}_c$ & M & $\infty$ & M \\
$K3$ (as $c = 6$ VOA) & M & $\infty$ & M \\
\bottomrule
\end{tabular}
\end{center}
For class~M seeds, the diagonal $T$-channel forces infinite depth in
$X_N^{\mathrm{id}}$. For class~G, L, or C seeds, invariant composite
fields in $X_N^{\mathrm{id}}$ may raise the fixed-point depth beyond
the seed value. A $G$-crossed completion may raise it again through
twist-field OPEs. These are two different questions.
\end{remark}


% ======================================================================
% Section: Koszulness
% ======================================================================

\section{Koszulness of the symmetric orbifold}
\label{sec:symn-koszulness}

\begin{proposition}[Tensor-product Koszulness and fixed-point descent;
\ClaimStatusConditional]\label{prop:symn-koszulness}
\index{symmetric orbifold!Koszulness}
Let $N\geq1$. If the seed algebra~$X$ is freely strongly generated
and hence chirally Koszul
(Proposition~\textup{\ref{prop:pbw-universality}}), then the
tensor product $X^{\otimes N}$ is chirally Koszul. The fixed-point
algebra $X_N^{\mathrm{id}}$ is chirally Koszul if the Reynolds
comparison
\[
 \barB(X_N^{\mathrm{id}})
 \longrightarrow
 \barB(X^{\otimes N})^{S_N}
\]
is a quasi-isomorphism.
\end{proposition}

\begin{proof}
The tensor product of freely strongly generated algebras is freely
strongly generated: if $\{a_i\}$ are free strong generators of~$X$,
then $\{a_i^{(r)}\}_{r=1}^N$ (the $i$-th generator in the $r$-th
copy) are free strong generators of $X^{\otimes N}$. PBW
universality (Proposition~\ref{prop:pbw-universality}) applies to
$X^{\otimes N}$, giving bar concentration. If the displayed Reynolds
comparison is a quasi-isomorphism, then
\[
 H^\bullet\barB(X_N^{\mathrm{id}})
 \cong H^\bullet(\barB(X^{\otimes N})^{S_N}).
\]
Over characteristic zero, $S_N$-invariants are exact, so the right
side is $H^\bullet\barB(X^{\otimes N})^{S_N}$ and remains
concentrated in the same bar degree.
\end{proof}

\begin{remark}[The two Koszul obstructions]
\label{rem:symn-full-koszulness}
Two independent obstructions remain after the tensor-product
calculation. First, fixed-point vertex algebras can have
invariant-theoretic relations even when the original tensor product
is freely strongly generated; this is measured by the Reynolds
comparison. Second, a $G$-crossed completion has twist fields and
binary twist OPEs, so its bar differential is not controlled by PBW
universality for $X^{\otimes N}$.
\end{remark}


% ======================================================================
% Section: The DMVV formula and second quantization
% ======================================================================

\section{The DMVV formula and second quantization}\label{sec:symn-dmvv}

The genus-$1$ elliptic genus of the full permutation-orbifold
conformal field theory has a closed expression in terms of the seed
elliptic genus. The DMVV product \cite{DMVV97} computes the
generating series over all~$N$ simultaneously:
for a seed theory~$X$ with elliptic genus
$Z_X(\tau, z) = \sum_{n,\ell} c_X(n,\ell)\, q^n y^\ell$,
\begin{equation}\label{eq:symn-dmvv}
 \sum_{N \geq 0} p^N\, Z(\operatorname{Sym}^N(X);\, \tau, z)
 \;=\; \prod_{\substack{n \geq 0,\; m \geq 1 \\ \ell}}
 \bigl(1 - q^n\, y^\ell\, p^m\bigr)^{-c_X(mn,\, \ell)}.
\end{equation}
The formula is a second-quantized lift: the genus-$1$ data of
the seed~$X$ (the Fourier coefficients $c_X(n, \ell)$) generate
the full orbifold partition function at all~$N$ through a
Borcherds-type infinite product. Its input and output are
elliptic genera, not ordered bar complexes.

\begin{remark}[DMVV as an elliptic-genus lift]
\label{rem:symn-dmvv-shadow}
\index{DMVV formula!shadow interpretation}
The logarithm of the DMVV product is the Hecke expansion of the
seed elliptic genus. It records the full genus-$1$ orbifold
spectrum, including commuting-pair twisted sectors. A shadow
coefficient of the ordered bar complex can be compared with this
series only after specifying a projection to the identity-sector
vacuum anomaly. There is no implication
$F_g(\operatorname{Sym}^N(X))=N F_g(X)$ for $g\geq2$ from DMVV.
\end{remark}

For $X = K3$, the full Jacobi lift is the reciprocal of the
unscaled Igusa cusp form~$\Phi_{10}^{\mathrm{un}}$. At the
Euler-characteristic specialization $z=0$ one obtains
$\sum_N p^N \chi(\operatorname{Hilb}^N(K3))
= \prod_{n \geq 1}(1 - p^n)^{-24}$.

\begin{proposition}[DMVV does not compute ordered-bar $\kappa$;
\ClaimStatusProvedHere]\label{prop:symn-dmvv-kappa}
\index{DMVV formula!kappa separation}
The coefficient of $p^N$ in the DMVV series~\eqref{eq:symn-dmvv}
is the full genus-$1$ orbifold elliptic genus. For $N>1$ it is not
the identity-sector ordered-bar scalar
$N\kappa(X)/24$. Any extraction of
Proposition~\ref{prop:symn-kappa} from DMVV requires an additional
projection to the diagonal vacuum-anomaly summand.
\end{proposition}

\begin{proof}
Taking the logarithm of~\eqref{eq:symn-dmvv} gives
\[
 \log\sum_{N\geq0}p^NZ(\operatorname{Sym}^N(X))
 =
 \sum_{m\geq1}\frac{p^m}{m}\,T_mZ_X .
\]
Thus the connected $p^N$ coefficient of the generating function is
$T_NZ_X/N$, while the coefficient of the original exponential also
contains products indexed by partitions of~$N$. The Hecke transform
includes non-identity covering data and twisted sectors. These
terms are part of the elliptic genus and have no canonical
identification with the ordered bar differential of
$X_N^{\mathrm{id}}$.
\end{proof}


% ======================================================================
% Section: Large-N limit and holography
% ======================================================================

\section{Large-$N$ limit and holography}\label{sec:symn-large-N}

The identity-sector scalar
$\kappa^{\mathrm{id}}_1(X_N^{\mathrm{id}})=N\kappa(X)$ grows
linearly with~$N$. The full symmetric-orbifold free energy is a
different object: it is extracted from the DMVV product and depends
on the polar spectrum and the choice of large-$N$ regime.

In a Cardy or holographic scaling limit one may write an expansion
\begin{equation}\label{eq:symn-large-N-expansion}
 \log Z(\operatorname{Sym}^N)
 \;=\; N f_0 + f_1 + N^{-1} f_2 + \cdots
\end{equation}
after fixing the ensemble and the normalization of $Z$. The
identity-sector scalar gives only the genus-$1$ diagonal anomaly:
\begin{equation}\label{eq:symn-genus-scaling}
 F^{\mathrm{id}}_1(\operatorname{Sym}^N(X))
 \;=\; \frac{N\kappa(X)}{24}.
\end{equation}
It does not determine the higher-genus coefficients $f_g$ of the
full orbifold partition function.

\begin{remark}[The $1/N$ expansion is not a bar theorem]
\label{rem:symn-genus-vs-1-over-N}
The string-theoretic rule $g_s\sim N^{-1/2}$ predicts genus-$g$
weights $N^{1-g}$ for connected closed-string amplitudes. The
ordered-bar scalar gives the boundary identity-sector anomaly
$N\kappa(X)$, not the whole closed-string genus expansion. A theorem
identifying higher shadow coefficients with bulk loop amplitudes
requires the Costello--Paquette bulk-boundary comparison and a
specified completion; it is not a consequence of DMVV or of
Proposition~\ref{prop:symn-kappa}.
\end{remark}

For $X = T^4$ or $K3$, Costello--Paquette
\cite{CP2020} identify the large-$N$ boundary
chiral algebra with the free symmetric orbifold, and the
bulk Kodaira--Spencer theory with the holographic dual. On the K3
identity lane, $\kappa^{\mathrm{id}}_1=2N$ while the
Brown--Henneaux central charge is $c=6N$; the anomaly ratio is
$\varrho=\kappa/c=1/3$, not the Virasoro ratio~$1/2$.


% ======================================================================
% Section: BTZ entropy and the shadow tower
% ======================================================================

\section{BTZ entropy and the shadow tower}\label{sec:symn-btz}

The BTZ black hole entropy
\begin{equation}\label{eq:symn-btz}
 S_{\mathrm{BH}}
 \;=\; 2\pi\sqrt{\frac{c \cdot n}{6}}
 \;=\; 2\pi\sqrt{\frac{\kappa \cdot n}{6\varrho}}
\end{equation}
where $n=L_0-c/24$ and $\varrho=\kappa/c$ on the chosen anomaly
lane. Cardy asymptotics determine the central charge. They determine
$\kappa$ only after the lane-specific ratio~$\varrho$ is fixed. For
the K3 symmetric-orbifold lane, $c=6N$, $\kappa=2N$, and
$\varrho=1/3$, giving
$S_{\mathrm{BH}}=2\pi\sqrt{Nn}$.

\begin{remark}[Logarithmic corrections and shadow data]
\label{rem:symn-btz-log-corrections}
\index{BTZ black hole!logarithmic corrections}
For a pure BTZ saddle with no extra light fields, the first
logarithmic correction has the standard form
\[
 S_{\mathrm{BH}}
 \;=\; 2\pi\sqrt{\frac{c n}{6}}
 \;-\; \frac{3}{2}\log\Bigl(\frac{c n}{6}\Bigr)
 \;+\; O(1),
\]
with the coefficient depending on the one-loop field content. Higher
ordered-bar shadows can supply boundary couplings for a bulk
comparison, but they do not universally determine the $O(1)$
Rademacher or one-loop terms of the full orbifold spectrum.
\end{remark}


% ======================================================================
% Section: The N=2 orbifold: explicit computations
% ======================================================================

\section{The \texorpdfstring{$N = 2$}{N = 2} orbifold:
explicit computations}\label{sec:symn-n2-explicit}

The simplest nontrivial symmetric orbifold is
$X_2^{\mathrm{id}}=(X\otimes X)^{\mathbb Z_2}$. The full
permutation-orbifold CFT also has the transposition sector $(12)$,
which carries one family of twist fields.

\begin{computation}[$\operatorname{Sym}^2(X)$ invariants]
\label{comp:sym2-explicit}
\index{symmetric orbifold!N=2 explicit}
Let $X$ be a chirally Koszul seed with central charge $c$ and
modular characteristic $\kappa(X)$. The fixed-point identity sector
has:
\begin{align}
 c(X_2^{\mathrm{id}})
 &= 2c(X), \label{eq:sym2-c} \\
 \kappa^{\mathrm{id}}_1(X_2^{\mathrm{id}})
 &= 2 \kappa(X), \label{eq:sym2-kappa} \\
 F^{\mathrm{id}}_1(\operatorname{Sym}^2(X))
 &= \kappa(X) / 12, \label{eq:sym2-f1}
\end{align}
by Proposition~\ref{prop:symn-kappa}. The transposition sector
contributes to the full orbifold genus-$1$ partition function, but
it does not enter the ordered bar complex of $X_2^{\mathrm{id}}$.
\end{computation}

\begin{remark}[The twist field at $N = 2$]
\label{rem:sym2-twist-field}
\index{twist field!N=2}
The transposition $(12) \in S_2$ introduces a twist field
$\sigma_2$ of conformal weight
\begin{equation}\label{eq:sym2-twist-weight}
 h(\sigma_2) \;=\; \frac{c(X)}{16}.
\end{equation}
The twist-field OPE is computed by the covering-space method:
the two-sheeted branched cover $\pi \colon \Sigma \to
\mathbb{P}^1$ ramified at the twist-field insertion points
maps the orbifold correlator to a correlator on the
covering surface~$\Sigma$. The leading
$\sigma_2$--$\sigma_2$ branch factor is
$(z-w)^{-c(X)/8}$, coming from the Casimir energy of the
two-sheeted cover. This is a branch exponent in the twisted module
category, not an ordinary meromorphic pole order in the fixed-point
chiral algebra.
\end{remark}

\begin{remark}[Covering space and the bar complex]
\label{rem:sym2-covering-bar}
\index{symmetric orbifold!covering space}
The covering-space technology has a bar-complex interpretation only
after the object has been chosen. For the fixed-point algebra,
\[
 \barB^{\mathrm{ord}}(X_2^{\mathrm{id}})
 = T^c\bigl(s^{-1}\overline{(X\otimes X)^{\mathbb Z_2}}\bigr),
\]
and the generators are invariant fields. There are no desuspended
twist fields in this complex.

For a $G$-crossed completion that adjoins $\sigma_2$, the bar
complex has three families of tensors:
\begin{enumerate}
\item invariant--invariant tensors inherited from
 $(X\otimes X)^{\mathbb Z_2}$;
\item The \emph{twist-twist} sector:
 tensor products of desuspended twist fields
 $s^{-1}\sigma_2 \otimes s^{-1}\sigma_2$, whose bar
 differential involves the $\sigma_2$-$\sigma_2$ OPE.
\item The \emph{mixed} sector: tensor products involving both
 untwisted and twist-field generators, with the bar
 differential involving mixed OPE data.
\end{enumerate}
The invariant-sector Koszul question is the Reynolds comparison of
Proposition~\ref{prop:symn-koszulness}. The completed-orbifold
Koszul question is the additional twist-OPE calculation.
\end{remark}


% ======================================================================
% Section: Hecke operators and the Fock space structure
% ======================================================================

\section{Hecke operators and the Fock space structure}
\label{sec:symn-hecke}

The DMVV formula~\eqref{eq:symn-dmvv} has a Hecke-operator form.
The generating series
$\sum_{N \geq 0} p^N Z(\operatorname{Sym}^N(X))$ can be
reorganized using the Hecke operator $T_m$ acting on the
seed elliptic genus:
\begin{equation}\label{eq:symn-hecke}
 \sum_{N \geq 0} p^N\, Z(\operatorname{Sym}^N(X);\, \tau, z)
 \;=\; \exp\Bigl(\sum_{m \geq 1} \frac{p^m}{m}\,
 T_m\, Z_X(\tau, z)\Bigr),
\end{equation}
where $T_m Z_X(\tau, z) = \sum_{ad = m,\, 0 \leq b < d}
Z_X\bigl(\frac{a\tau + b}{d},\, az\bigr)$ is the $m$-th
Hecke operator acting on the Jacobi form $Z_X$.

\begin{remark}[Fock space interpretation]
\label{rem:symn-fock}
\index{symmetric orbifold!Fock space}
The exponential in~\eqref{eq:symn-hecke} identifies the
symmetric orbifold partition function with the partition
function of a \emph{bosonic Fock space} built from the
Hecke-transformed seed elliptic genus. The ``single-particle''
object is the graded vector space counted by the seed elliptic
genus, not the ordered bar complex of the seed VOA. The Hecke
operators organize covering data of the torus. They do not act as
creation operators on the chiral bar complex and do not imply
$S_r(\operatorname{Sym}^N(X))=N S_r(X)$.
\end{remark}

\begin{remark}[K3 seed and the genus-$2$ Igusa character]
\label{rem:symn-k3-igusa-character}
For the K3 seed one has the classical DMVV identity
\[
 \sum_{N\geq 0} p^N\,Z\bigl(\operatorname{Sym}^N(K3);\tau,z\bigr)
 \;=\;
 \frac{1}{\Phi_{10}^{\mathrm{un}}(\tau,z,\sigma)}.
\]
Two boundary regimes of this same Siegel form must be kept distinct. The
Fourier--Jacobi cusp $p\to 0$ gives
\[
 \frac{1}{\Phi_{10}^{\mathrm{un}}(\tau,z,\sigma)}
 \;=\;
 \frac{p^{-1}}{\phi_{10,1}(\tau,z)} + O(p^0),
 \qquad
 \phi_{10,1}(\tau,z)=\eta(\tau)^{18}\vartheta_1(\tau,z)^2.
\]
The separating divisor $z\to 0$ gives
\[
 \frac{1}{\Phi_{10}^{\mathrm{un}}(\tau,z,\sigma)}
 \;=\;
 -\frac{1}{4\pi^2 z^2\,\eta(\tau)^{24}\eta(\sigma)^{24}} + O(z^0).
\]
Volume~III uses the same form as the genus-$2$ character of the K3 chiral
bialgebra $\mathbf{H}_{\Delta_5}$; the first limit isolates the Jacobi
coefficient, while the second is the genuine separating degeneration.
\end{remark}

\begin{proposition}[Hecke operators and the identity scalar;
\ClaimStatusProvedHere]\label{prop:symn-hecke-kappa}
\index{Hecke operator!identity scalar}
The Hecke operator $T_1 = \mathrm{id}$ reproduces the
seed contribution, so the $p^1$ term of~\eqref{eq:symn-hecke}
is $Z_X(\tau, z)$ and the connected genus-$1$ amplitude at
$N = 1$ has identity scalar $F^{\mathrm{id}}_1(X)=\kappa(X)/24$.
For $N>1$, the Hecke coefficient is not the ordered-bar scalar
$N\kappa(X)/24$ without the identity-sector projection of
Proposition~\ref{prop:symn-kappa}.
\end{proposition}

\begin{proof}
The logarithm of~\eqref{eq:symn-hecke} gives the connected
generating function:
$\sum_{N \geq 1} p^N F_1^{\mathrm{conn}}
(\operatorname{Sym}^N(X))
= \sum_{m \geq 1} (p^m / m) T_m Z_X$.
At order $p^N$: the connected genus-$1$ amplitude receives
contributions from the $m = N$ Hecke operator acting on the
seed elliptic genus. Its constant terms are divisor-weighted
Hecke sums, not the diagonal copy sum in the ordered bar complex.
The ordered-bar scalar is recovered only by isolating the diagonal
identity-sector anomaly before taking the full orbifold trace.
\end{proof}


% ======================================================================
% Section: The bar complex of Sym^N(X)
% ======================================================================

\section{The bar complex of
\texorpdfstring{$\operatorname{Sym}^N(X)$}{Sym\^{}N(X)}}
\label{sec:symn-bar-complex}

The ordered bar complex depends on which orbifold object is meant.
For the fixed-point chiral algebra it is the bar complex of
$X_N^{\mathrm{id}}=(X^{\otimes N})^{S_N}$; for a $G$-crossed
completion it is a larger complex with twisted generators.

\begin{computation}[Fixed-point bar complex]
\label{comp:symn-bar-decomposition}
\index{bar complex!symmetric orbifold}
The fixed-point ordered bar complex is
\[
 \barB^{\mathrm{ord}}(X_N^{\mathrm{id}})
 =
 T^c\bigl(s^{-1}\overline{(X^{\otimes N})^{S_N}}\bigr).
\]
Its generators are invariant combinations of tensor-product fields,
for example diagonal generators
$a_\Delta=\sum_i a^{(i)}$ and higher invariant normally ordered
polynomials. The natural comparison map
\[
 \barB^{\mathrm{ord}}(X_N^{\mathrm{id}})
 \longrightarrow
 \barB^{\mathrm{ord}}(X^{\otimes N})^{S_N}
\]
is the Reynolds descent map of
Proposition~\ref{prop:symn-koszulness}. Twist fields
$\sigma_\lambda$ are absent from this complex. They enter only after
a $G$-crossed completion has been chosen.
\end{computation}

\begin{remark}[Collision residues of the fixed-point and twisted objects]
\label{rem:symn-r-matrix}
\index{symmetric orbifold!$R$-matrix}
The collision residue $r^{\mathrm{coll}}(z) =
\operatorname{Res}^{\mathrm{coll}}_{0,2}
(\Theta_{X_N^{\mathrm{id}}})$ has a diagonal fixed-point channel.
For the stress tensor,
\[
 T_\Delta(z)T_\Delta(w)
 \sim
 \frac{N\,c(X)/2}{(z-w)^4}
 + \frac{2T_\Delta(w)}{(z-w)^2}
 + \frac{\partial T_\Delta(w)}{z-w}.
\]
Thus the Virasoro collision residue has the same pole order as the
seed and scalar coefficient multiplied by~$N$.

Twist-twist channels are not meromorphic channels of
$X_N^{\mathrm{id}}$. In a completed orbifold category, the
$N=2$ field satisfies
\[
 T(z)\sigma_2(w)\sim
 \frac{(c/16)\sigma_2(w)}{(z-w)^2}
 +\frac{\partial\sigma_2(w)}{z-w},
 \qquad
 \sigma_2(z)\sigma_2(w)\sim
 (z-w)^{-c/8}(1+\cdots).
\]
The second formula is a branch expansion on the cover, not a pole of
integer order in the fixed-point VOA.
\end{remark}

\begin{remark}[Koszulness and the permutation orbifold
functor]\label{rem:symn-orbifold-functor}
\index{symmetric orbifold!orbifold functor}
The fixed-point construction defines the functor
$X\mapsto (X^{\otimes N})^{S_N}$ on vertex operator algebras. Its
Koszul question is whether the Reynolds comparison for bar complexes
is a quasi-isomorphism. A full permutation-orbifold CFT is not this
functor; it is the fixed-point algebra together with twisted modules,
and a $G$-crossed completion must specify their OPEs.

At the level of identity-sector modular characteristics,
$\kappa^{\mathrm{id}}_1(X_N^{\mathrm{id}})=N\kappa(X)$ by
Proposition~\ref{prop:symn-kappa}. Higher shadow coefficients and
twisted sectors are not functorial consequences of that scalar
identity.
\end{remark}

\begin{remark}[Comparison with the wreath product]
\label{rem:symn-wreath}
\index{symmetric orbifold!wreath product}
The permutation action is the vertex-algebra analogue of the wreath
product $\Gamma \wr S_N$ in finite group theory: the group $\Gamma$
(the automorphism group of~$X$) is wrapped by the symmetric group
$S_N$. When $\Gamma$ is trivial, the fixed-point algebra is
$(X^{\otimes N})^{S_N}$ and the full CFT still has permutation
twisted modules. When
$\Gamma = \mathbb{Z}_2$ (e.g., $X = \beta\gamma$ with its
$\mathbb{Z}_2$ symmetry), the wreath product
$\mathbb{Z}_2 \wr S_N$ acts on the tensor product, and the
fixed-point algebra is $(X^{\otimes N})^{\mathbb Z_2\wr S_N}$.
The completed CFT has both copywise $\mathbb Z_2$ twisted sectors
and permutation twisted sectors. Its Koszulness is a problem about
that completed OPE algebra, not a corollary of the fixed-point
scalar identity.
\end{remark}


% ======================================================================
% Section: Open problems
% ======================================================================

\section{Open problems}\label{sec:symn-open-problems}

\begin{enumerate}
\item \textbf{Reynolds descent.}
Decide when
$\barB((X^{\otimes N})^{S_N})\to\barB(X^{\otimes N})^{S_N}$
is a quasi-isomorphism. This is the fixed-point Koszulness
obstruction.

\item \textbf{Twisted-sector Koszul dual.}
Construct the $G$-crossed bar complex when twist fields are adjoined
and compute its Koszul dual from the twist-field OPE and
covering-space structure.

\item \textbf{Higher-genus orbifold amplitudes.}
The DMVV formula is a genus-$1$ elliptic-genus statement. Determine
which higher-genus amplitudes of the full orbifold CFT are fixed by
seed data and which require new twisted-sector sewing input.

\item \textbf{Exact shadow depth for non-M seeds.}
For class~G, L, or C seeds, determine whether invariant composite
fields in $X_N^{\mathrm{id}}$ increase the fixed-point shadow depth.
For completed orbifolds, repeat the calculation with twist-field
channels.
\end{enumerate}
